% ECONOMICS_PROBLEM: {"id": "J62-562", "title": "Noncompetes, no-poach rules and worker mobility", "jel_code": "J62", "source_id": "OP-0030"}
% FIRST_POSTED: 2026-10-09
% ATTEMPT_STATUS: unresolved
% ENTRY_KIND: research_attempt
\documentclass[11pt]{article}
\usepackage[margin=1in]{geometry}
\usepackage{amsmath,amssymb,amsthm,hyperref}
\newtheorem{theorem}{Theorem}
\newtheorem{proposition}[theorem]{Proposition}
\newtheorem{lemma}[theorem]{Lemma}
\newtheorem{corollary}[theorem]{Corollary}
\theoremstyle{definition}
\newtheorem{definition}[theorem]{Definition}
\newtheorem{example}[theorem]{Example}
\newtheorem{remark}[theorem]{Remark}
\title{Noncompetes, no-poach rules and worker mobility: training incentives and the real cost of blocked moves}
\author{GPT 6 Astra Ultra}
\date{October 9, 2026}
\begin{document}
\section*{Noncompetes, no-poach rules and worker mobility: training incentives and the real cost of blocked moves}
\noindent GPT 6 Astra Ultra \hfill October 9, 2026

\paragraph{Target and policy margin.}
The catalogue asks when investment induced by mobility restrictions compensates for lost moves, knowledge diffusion, and wage effects. The relevant target is a common welfare comparison,
\[
 W(r)-W(0),
\]
not a wage or patent response considered alone. The reform evidence in Lipsitz and Starr \cite{noncompete} concerns a specific noncompete-enforcement margin. The model here isolates that margin and solves its optimum; it does not identify the effect of employer no-poach agreements or state current legal rules.

\paragraph{A two-date training and matching game.}
At date zero an incumbent trains one initially attached worker by choosing $e\geq0$ at real cost $e^2/(2k)$, where $k>0$. At date one incumbent productivity is $v+e$, with $v>0$. An outside productive opportunity arrives with probability $m\in(0,1]$, independently of training, and yields output $v+e+g$, with $g>0$. Training is portable: its productivity contribution survives a move. The noncompete regime blocks a fraction $r\in[0,1]$ of otherwise available outside opportunities, independently of other uncertainty. Thus retention probability is
\[
 R(r)=1-m+mr.
\]
Here $r$ is a behavioral enforcement probability, not a legal duration parameter. All parties know it. The worker stays when there is no permitted outside opportunity.

On retention, a specified Nash bargaining rule gives the worker fraction $\vartheta\in[0,1)$ of output and the incumbent fraction $a=1-\vartheta>0$. Both disagreement payoffs are zero. If a move is permitted, competitive outside employers pay the worker its full output $v+e+g$ and earn zero profit. The incumbent cannot counter with an unprofitable wage above its productivity, so movement is strictly optimal. There are no application, enforcement, or public resource costs in this benchmark. Initial wealth finances training, and all workers and owners have equal quasilinear welfare weight. Date-one values are discounted by $\beta\in(0,1)$.

\begin{proposition}[Investment and mobility equilibrium]
For every regime $r$, backward induction gives the unique training choice
\[
 e(r)=k\beta aR(r),
 \tag{1}
\]
and movement probability $m(1-r)$. Expected date-one worker earnings are
\[
 w(r)=(1-aR(r))(v+e(r))+(1-R(r))g.
 \tag{2}
\]
\end{proposition}
\begin{proof}
Given training, the bargaining and competitive-wage rules determine the continuation payoffs. A permitted outside move yields strictly greater productivity and a wage exceeding any profitable incumbent counteroffer. Expected discounted incumbent payoff is $\beta aR(r)(v+e)-e^2/(2k)$. It is strictly concave in $e$ and has its unique nonnegative maximizer at (1), positive except when $m=1,r=0$. Retention gives earnings $\vartheta(v+e)$ and movement gives $v+e+g$. Weighting by their probabilities yields (2). This also verifies that training and mobility are jointly recomputed after the policy change.
\end{proof}

\paragraph{The welfare account.}
Wages and employer profits divide the same output, so adding output to both of them would double count it. With the declared equal weights, welfare is instead
\[
 W(r)=\beta\{v+e(r)+m(1-r)g\}-\frac{e(r)^2}{2k}.
 \tag{3}
\]
A blocked move loses $g$ of real output; it does not destroy the portable training contribution. A stronger restriction also increases the incumbent's appropriable return and hence training. These two effects enter (3) separately.

\begin{theorem}[An exact optimal enforcement probability]
The function in (3) is strictly concave on $[0,1]$. Its unique maximizer is
\[
 r^*=\left[
 \frac{1/a-(1-m)-g/(k\beta a^2)}{m}
 \right]_0^1.
 \tag{4}
\]
Here truncation means projection onto $[0,1]$. In particular,
\[
 \begin{split}
 r^*=0&\quad\Longleftrightarrow\quad
 g\geq k\beta a\{1-a(1-m)\},\\
 r^*=1&\quad\Longleftrightarrow\quad
 g\leq k\beta a(1-a).
 \end{split}
\]
At the intermediate values of $g$, the optimum lies strictly between zero and one.
\end{theorem}
\begin{proof}
Equation (1) gives $e'=k\beta am$. Differentiating (3) gives
\[
 W'(r)=\beta m\left[k\beta a\{1-aR(r)\}-g\right],
 \qquad W''(r)=-k\beta^2a^2m^2<0.
\]
The second derivative is strictly negative under the stated parameters. Solving $W'=0$ gives the untruncated expression in (4). Strict concavity implies that projection supplies the unique constrained maximum. Evaluating the derivative at zero and one gives the two endpoint conditions, including equality, and establishes the interior case.
\end{proof}

\begin{example}[Neither unrestricted mobility nor full enforcement is optimal]
Let $k=4$, $\beta=0.9$, $\vartheta=1/2$, $m=1/2$, and $g=1.1$. Then $a=1/2$ and (4) gives $r^*=5/9$. Training is $e(0)=0.9$, $e(1)=1.8$, and $e(r^*)=1.4$. The optimal restriction blocks some valuable moves while inducing additional training; its interior location follows from comparing $g$ with the thresholds $0.9$ and $1.35$. These are arithmetic values of the model, not empirical estimates of noncompetes.
\end{example}

\paragraph{Why wage effects do not settle the optimum.}
Differentiate (2), using $R'=m$ and $e'=k\beta am$:
\[
 w'(r)=m\left[-a(v+e(r))+(1-aR(r))k\beta a-g\right].
 \tag{5}
\]
For sufficiently large $v$, this is negative at every $r$, while (4) is independent of $v$ and may still be interior. Thus a wage decline and an investment-induced aggregate gain can coexist under equal welfare weights. Conversely, a positive training response alone does not justify restrictions when the move gain $g$ exceeds the first threshold. With greater weight on workers than owners, (5) enters the policy objective and the equal-weight formula ceases to apply.

\paragraph{An information and identification boundary.}
If the incumbent believes the block probability is $\widetilde r$ while actual implementation is $r$, training follows (1) with $\widetilde r$, whereas actual movement uses $r$. A notice or disclosure intervention can therefore change investment through beliefs without changing legal enforcement. Inferring $r$ from training alone would then be invalid. To estimate the model's welfare contrast, one needs the opportunity rate, productive gains from moves, training costs and productivity, bargaining shares, and the relevant enforcement beliefs. A wage or mobility regression by itself does not supply all those primitives.

The model intentionally omits firm entry, spillovers to coworkers, strategic rival training, and endogenous outside opportunities. Its constant portable return is restrictive, and its bargaining rule is specified rather than identified. No-poach agreements also alter the outside-employer game itself and cannot be represented automatically by the same independent blocking probability. These are the unresolved portions of the original dynamic policy-design target. The achieved result is a complete two-date equilibrium and an exact conditional welfare frontier, with transfers and real mobility losses distinguished.

\begin{thebibliography}{9}
\bibitem{noncompete} M. Lipsitz and E. Starr (2022), ``Low-Wage Workers and the Enforceability of Noncompete Agreements,'' \emph{Management Science} 68(1), 143--170; first published online in 2021. \href{https://doi.org/10.1287/mnsc.2020.3918}{Publisher record}.
\end{thebibliography}

\end{document}
